% !TEX program = lualatex
% !TeX encoding = UTF-8
\PassOptionsToPackage{unicode}{hyperref}
\PassOptionsToPackage{bookmarksnumbered,bookmarksopen}{hyperref}
\documentclass[12pt,a4paper]{article}

% ============================================================
% 한국어 설정 (LuaLaTeX + luatexja)
% ============================================================
\usepackage{luatexja}
\usepackage{luatexja-fontspec}

% 한국어 고딕체 (sans) – Noto Sans KR
\setsansjfont{Noto Sans KR}[
UprightFont = * Regular,
BoldFont = * Bold,
Script = Hangul,
Language = Korean
]

% 한국어 명조체 (serif) – Noto Serif KR
\IfFontExistsTF{Noto Serif KR}{
	\setmainjfont{Noto Serif KR}[
	UprightFont = * Regular,
	BoldFont = * Bold,
	Script = Hangul,
	Language = Korean
	]
}{
	% fallback
	\setmainjfont{Noto Sans KR}[
	UprightFont = * Regular,
	BoldFont = * Bold,
	Script = Hangul,
	Language = Korean
	]
}

% 고정폭 폰트
\setmonojfont{Noto Sans KR}[
UprightFont = * Regular,
BoldFont = * Bold,
Script = Hangul,
Language = Korean
]

% 라틴 문자 폰트
\setmainfont{Times New Roman}[Ligatures=TeX]
\setsansfont{Arial}[Ligatures=TeX]
\setmonofont{Latin Modern Mono}[
WordSpace={1,0,0},
HyphenChar=None,
PunctuationSpace=WordSpace
]

% ============================================================
% 기타 패키지
% ============================================================
\usepackage{amsmath,amssymb,amsthm,mathtools}
\usepackage{geometry}
\usepackage{booktabs}
\usepackage{hyperref}
\usepackage{url}
\usepackage{xcolor}
\usepackage{enumitem}
\usepackage{tikz}
\usetikzlibrary{decorations.pathreplacing,arrows.meta,positioning,shapes.geometric}
\usepackage{float}

\geometry{a4paper, margin=1in}

% ============================================================
% YUCT 매크로
% ============================================================
\newcommand{\Keff}{K_{\mathrm{eff}}}
\newcommand{\kappac}{\kappa_c}
\newcommand{\eps}{\varepsilon}
\newcommand{\df}{d_{\!f}}

% ============================================================
% 정리 환경 (한국어)
% ============================================================
\newtheorem{theorem}{정리}[section]
\newtheorem{lemma}[theorem]{보조정리}
\newtheorem{definition}[theorem]{정의}
\newtheorem{corollary}[theorem]{따름정리}
\newtheorem{proposition}[theorem]{명제}
\theoremstyle{remark}
\newtheorem{remark}[theorem]{주석}
\newtheorem{example}[theorem]{예시}

% ============================================================
% 하이퍼링크 설정
% ============================================================
\hypersetup{colorlinks=true,linkcolor=blue,citecolor=blue,urlcolor=blue}

\begin{document}
	
	\begin{titlepage}
		\centering
		\vspace*{1cm}
		{\Huge\bfseries 부록 S2: 2차원 쇼트키 헤테로구조에서 조정의 프랙탈 차원과 \\ YUCT 스케일링 지수 $\beta = 2/3$ ($\approx 0.67$)의 예측}\\[1cm]
		{\large Alexey V. Yakushev}\\[0.3cm]
		{\normalsize Yakushev Research, \url{https://yuct.org/}}\\[1.5cm]
		{\normalsize DOI: \texttt{10.5281/zenodo.18444598}}\\[0.5cm]
		{\normalsize 2026년 5월\\버전 1.0}\\[2cm]
		
		\begin{minipage}{0.9\textwidth}
			\begin{abstract}
				\noindent
				Ang, Yang \& Ang (Phys. Rev. Lett. 121, 056802, 2018)은
				2차원 쇼트키 헤테로구조에서 보편적인 전류-온도 스케일링 법칙
				$\log(\mathcal{I}/T^{3/2}) \propto -1/T$ (측면 접촉) 및
				$\log(\mathcal{I}/T^{1}) \propto -1/T$ (운동량 비보존 수직 접촉)을
				발견했다. 본 부록에서는 두 지수가 수송을 제어하는 조정 네트워크의
				유효 프랙탈 차원 $d_f$에 의해 결정됨을 보여준다.
				YUCT 보편 오차 스케일링 법칙 $\varepsilon = \kappa_c\alpha
				K_{\mathrm{eff}}^{-\beta}$ ($\beta = 2/3 \approx 0.67$)은
				$\beta = 2/d_f$를 통해 $d_f$와 연결된다. 실험적 전류 스케일링으로부터
				측면 접촉에서는 $d_f = 3$, 수직 접촉에서는 $d_f = 2$를 추출하며,
				각각 $\beta = 2/3$ 및 $\beta = 1$을 제공한다. 이 분석은
				측면 쇼트키 다이오드에서 상대 전류 변동이
				$\varepsilon_{\mathrm{fluc}} \propto K_{\mathrm{eff}}^{-2/3}$을
				따를 것이라고 예측하며, 이는 표준 노이즈 분광법으로 검증 가능한
				반증 가능한 예측이다. 본 부록은 2차원 쇼트키 헤테로구조를
				보편적인 YUCT 지수 $\beta = 2/3$ ($0.67$)을 나타내거나 나타낼 것으로
				예측되는 계의 패밀리에 포함시켜, 이론의 적용 범위를 고체 나노전자공학
				영역으로 확장한다.
			\end{abstract}
			\vspace{1cm}
			\noindent\small\textbf{키워드:} YUCT, 조정 효율, 보편 스케일링,
			$\beta = 2/3$, 0.67, 쇼트키 헤테로구조, 2차원 물질,
			프랙탈 차원, KPZ 보편성.
		\end{minipage}
	\end{titlepage}
	
	\tableofcontents
	\newpage
	
	\section{서론}
	Yakushev 통합 조정 이론 (YUCT)은 임의의 조정 과정이 보편 오차 스케일링 법칙을
	따른다고 가정한다:
	\begin{equation}
		\varepsilon = \kappa_c \, \alpha \, \Keff^{-\beta},
		\qquad \beta = 2/3 \approx 0.67,\;\; \kappa_c = 1/3,
		\label{eq:yuct-law}
	\end{equation}
	여기서 $\Keff$는 조정 효율, $\alpha$는 시스템 특정 전인자이다.
	이 법칙은 화학 결합에서 우주론적 변동에 이르기까지 다양한 시스템에서
	검증되었다 (부록~L,~M,~E,~F,~G,~V).
	
	Ang, Yang and Ang~\cite{Ang2018}은 2차원 물질 기반 쇼트키 헤테로구조에
	대한 보편적인 전류-온도 스케일링 법칙을 유도했다:
	\begin{align}
		\text{측면 접촉:}&\quad \log\Bigl(\frac{\mathcal{I}}{T^{3/2}}\Bigr) \propto -\frac{1}{T},
		\label{eq:ang-lat}\\
		\text{수직 접촉 ($k_\parallel$ 비보존):}&\quad \log\Bigl(\frac{\mathcal{I}}{T^{1}}\Bigr) \propto -\frac{1}{T}.
		\label{eq:ang-vert}
	\end{align}
	지수 $\beta_{\mathrm{IV}} = 3/2$ 및 $1$은 특정 2차원 물질에 의존하지 않으며,
	전류 경로의 기하학을 반영한다.
	
	본 부록에서는 지수 $\beta_{\mathrm{IV}}$가 독립적인 수가 아니라,
	수송을 제어하는 조정 네트워크의 유효 프랙탈 차원 $d_f$에 의해
	결정됨을 보여준다. 실험값 $\beta_{\mathrm{IV}}$로부터 $d_f$를 추출하고,
	$d_f$로부터 조정 오차 지수 $\beta_{\mathrm{err}} = 2/d_f$를 예측한다.
	측면 접촉에서는 $\beta_{\mathrm{err}} = 2/3$이 되며, 이는 정확히
	YUCT 보편값이다. 이 예측으로부터 도출되는 변동 스케일링은 이론의
	반증 가능한 테스트를 구성한다.
	
	\section{쇼트키 장벽을 위한 YPSDC 사전}
	YPSDC (동기 분산 조정의 Yakushev 프로토콜)의 언어에서, 측면 쇼트키
	헤테로구조는 다음 구성 요소를 가진 조정 시스템이다:
	\begin{itemize}
		\item \textbf{사전 $\mathcal{D}$:} 2차원 평면 내의 일전자 상태들의 집합으로,
		에너지 $\epsilon > \Phi_B$이고 파동 벡터 $\mathbf{k}$가
		면내 분산 (\ref{eq:dispersion})을 만족하는 것.
		\item \textbf{색인 $\kappa$:} 전자를 사전으로 여기시키는 열에너지 $k_B T$.
		\item \textbf{공명 $R$:} 장벽 투과 중 접선 운동량 $k_y$의 보존.
	\end{itemize}
	조정된 행동은 에너지 조건 $\epsilon > \Phi_B$와 운동량 조건 $k_x > 0$ (장벽 방향)
	을 동시에 만족하는 전자의 열전자 방출이다. 이 조정의 효율이 역방향 포화 전류를
	결정한다.
	
	\section{전류의 실험적 스케일링과 조정 효율}
	
	일반적인 등방성 2차원 분산
	\begin{equation}
		\epsilon(\mathbf{k}) = \sum_n c_n |\mathbf{k}|^n,
		\label{eq:dispersion}
	\end{equation}
	에 대해, Ang 등은 활성 채널 수 (전류에 기여할 수 있는 것)가
	보편적으로
	\begin{equation}
		N_{\mathrm{ch}} \propto T^{1/2}.
		\label{eq:Nch}
	\end{equation}
	와 같이 스케일함을 보였다. 장벽 위의 열적으로 접근 가능한 전체 상태 수는
	$T$에 비례한다 (볼츠만 꼬리로부터). 따라서 성공적으로 조정된
	상태의 비율
	\begin{equation}
		\Keff(T) = \frac{N_{\mathrm{ch}}}{N_{\mathrm{tot}}} \propto T^{-1/2},
		\label{eq:keff-T}
	\end{equation}
	이 이 과정의 조정 효율로 기능한다. 측정되는 전류는 다음 형태를 취한다:
	\begin{equation}
		\mathcal{I} = \mathcal{I}_0 \, \Keff(T)^{-3} \, e^{-\Phi_B/k_B T},
		\label{eq:I-via-keff}
	\end{equation}
	이는 실험 법칙 (\ref{eq:ang-lat})을 재현한다.
	
	노이즈 측정에 적합한 $\Keff$의 운영적 정의는
	\begin{equation}
		\Keff = \frac{I}{I_0},
		\label{eq:keff-operational}
	\end{equation}
	여기서 $I$는 측정 전류, $I_0$는 조정 제한이 없었다면 흐를 참조 전류이다.
	$I_0$는 수직 접촉의 고온 거동 ($\Keff(T) \propto T^{-0}$ 점근적)을
	측정 온도로 외삽하거나, 동일한 장벽 높이를 가진 3차원 쇼트키 다이오드를
	이용한 독립적 교정으로부터 얻을 수 있다. 이 운영적 정의는 $\Keff$가
	이론적 스케일링 (\ref{eq:keff-T})에 의존하지 않고 전류-전압 특성으로부터
	직접 결정될 수 있게 한다.
	
	\section{조정 네트워크의 프랙탈 차원}
	
	YUCT (부록~L)의 중심 결과는 오차 스케일링 지수가 조정 네트워크의
	프랙탈 차원 $d_f$에 의해 결정된다는 것이다:
	\begin{equation}
		\beta_{\mathrm{err}} = \frac{2}{d_f}.
		\label{eq:beta-df}
	\end{equation}
	열적 창 내의 활성 조정 채널의 밀도는 네트워크의 기하학적 특징을
	담는다. 등방성 $d_f$ 차원 공간에서 온도 $T$에서 접근 가능한 상태 밀도는
	$T^{d_f/2 - 1}$로 스케일한다. 측면 쇼트키 장벽에서는,
	운동량 보존 ($\cos\theta$ 인자)이라는 추가적 운동학적 제약이
	유효 지수를 1만큼 감소시키므로, 운동 지수는
	\begin{equation}
		\beta_{\mathrm{IV}} = \frac{d_f}{2}.
		\label{eq:bIV-df}
	\end{equation}
	\textbf{식 (\ref{eq:bIV-df})는 전류 경로와 프랙탈 차원 $d_f$의
		조정 네트워크 간의 기하학적 유사성에 의해 동기부여된 가설이다.
		YUCT 작용으로부터의 엄밀한 유도는 아직 수행되지 않았다; 따라서 이 관계는
		검증 가능한 예측 (섹션~\ref{sec:predictions})을 이끌고 미래의 이론적
		연구를 촉구하는 추측으로 간주되어야 한다.}
	
	식 (\ref{eq:beta-df})과 (\ref{eq:bIV-df})을 결합하여 기본 관계를 얻는다:
	\begin{equation}
		\boxed{\beta_{\mathrm{IV}} \cdot \beta_{\mathrm{err}} = 1}.
		\label{eq:relation}
	\end{equation}
	
	실험값 $\beta_{\mathrm{IV}} = 3/2$로부터 다음을 얻는다:
	\begin{equation}
		d_f = 2\beta_{\mathrm{IV}} = 3, \qquad
		\beta_{\mathrm{err}} = \frac{2}{d_f} = \frac{2}{3} \approx 0.67.
	\end{equation}
	따라서 측정된 전류 스케일링은 유효 프랙탈 차원 $d_f = 3$의 조정 네트워크를
	결정하며, 이는 이제 임의의 조정 오차의 척도 (예: 상대 전류 변동)가
	YUCT의 거듭제곱 법칙에 지수 $2/3$로 따라야 함을 의미한다. 이는 \textbf{예측}이며,
	사후 예측이 아니다. 왜냐하면 이 시스템들에서는 변동 측정이 아직 수행되지
	않았기 때문이다.
	
	$k_\parallel$이 비보존인 수직 접촉에서는 각도 제한이 사라지고,
	$\beta_{\mathrm{IV}} = 1$, 따라서 $d_f = 2$, $\beta_{\mathrm{err}} = 1$이 된다.
	YUCT는 그러한 접촉에서의 전류 변동이 $\Keff^{-1}$로 스케일하고,
	$2/3$ 법칙과 구별된다고 예측한다.
	
	\section{전류 변동의 예측}
	\label{sec:predictions}
	조정 효율 (\ref{eq:keff-T})에 적용된 YUCT 오차 법칙 (\ref{eq:yuct-law})은
	측면 쇼트키 다이오드에서 저주파 전류 노이즈 전력 $S_I$에 대한
	정량적 예측을 제공한다:
	\begin{equation}
		\frac{S_I}{I^2} \propto \Keff^{-2/3} \propto T^{1/3}.
		\label{eq:noise-pred}
	\end{equation}
	식 (\ref{eq:keff-operational})의 운영적 정의 $\Keff = I/I_0$를 사용하면,
	노이즈 스케일링은 다음과 같이 다시 쓰여질 수 있다:
	\begin{equation}
		\frac{S_I}{I^2} = C \, \left(\frac{I}{I_0}\right)^{-2/3} + \text{const},
		\label{eq:noise-operational}
	\end{equation}
	여기서 $C$는 재료 특정 상수로 1 정도, 가산 상수는 기약 열잡음을 고려한다.
	이 형태는 표준 노이즈 분광 실험에서 직접 접근 가능하다:
	$I$와 $I_0$는 $I$--$V$ 특성으로부터 측정되고, $S_I$는 일정 온도에서
	바이어스의 함수로 기록된다. 예측은 $S_I/I^2$를 $(I/I_0)^{-2/3}$에 대해
	플롯했을 때 선형이 되고, 그 기울기는 온도에 의존하지 않는다는 것이다.
	이는 지수의 피팅을 전혀 필요로 하지 않는 엄격한 반증 가능한 테스트이다.
	
	\section{고찰: 스케일링 지수의 통일}
	
	$3/2$ 지수의 출현은 쇼트키 다이오드에 특유한 것이 아니라 더 깊은 조직 원리를
	나타내며, 이는 Kardar–Parisi–Zhang (KPZ) 표면 성장의 보편성 부류에
	독립적으로 나타난다는 사실에서도 분명하다~\cite{KPZ}. KPZ에서는 동적 지수 $z = 3/2$가
	성장 계면의 프랙탈 거칠기를 지배한다. YUCT는 이러한 현상들에 대해 통일된 해석을 제공한다.
	\textbf{성장 전선은 $d_f = 3$ 차원의 조정 네트워크로 개념화될 수 있으며,
		여기서 운동 지수 $\beta_{\mathrm{IV}} = d_f/2 = 3/2$가 우리의 추측 관계 (\ref{eq:bIV-df})
		로부터 자연스럽게 나타난다.} 따라서 KPZ 지수와 쇼트키 지수는 단순히 유사한 것이 아니라,
	완전히 다른 물리적 실현에서 동일한 YUCT 오차 법칙 $\varepsilon \propto \Keff^{-2/3}$의
	두 가지 운동적 서명이다. \textbf{이 연결은 현재 가설이며, YUCT 라그랑지안을 KPZ 방정식에
		직접 연결하는 추가 이론적 작업을 필요로 한다.}
	
	KPZ 예측을 더 구체적으로 만들기 위해, $\Keff$를 성장 계면에서 포화 상관 모드의
	비율과 명시적으로 동일시할 수 있다. 측면 크기 $L$의 KPZ 계면의 포화 영역에서,
	폭 $W(L,t)$는 정상 상태값 $W_{\text{sat}}(L) \simeq \sqrt{\langle h^2 \rangle} \propto L^{\alpha}$
	($\alpha = 1/2$)에 접근한다. 변동 진폭 $\delta W = W(t) - W_{\text{sat}}$는
	활성 (비포화) 모드의 수에 의해 제어된다. 우리는 $\Keff$를
	$1 - \nu_{\text{unsat}}/\nu_{\text{tot}}$, 즉 상관된 정상 상에 들어간 모드의 비율로
	동일시할 것을 제안한다. 그러면 YUCT는
	\begin{equation}
		\frac{\delta W}{W_{\text{sat}}} \propto \Keff^{-2/3},
		\label{eq:kpz-prediction}
	\end{equation}
	를 예측한다. 여기서 $\Keff$는 폭의 시간 감쇠로부터 $\Keff \approx 1 - (dW/dt)_{\text{norm}}$
	로 추정될 수 있다. 이 예측은 수치 시뮬레이션 및 동적 거칠기의 고분해능 실험
	(예: 전착, 분자선 에피택시)에서 검증 가능하다.
	
	이 통일은 그림~\ref{fig:hierarchy}에 나타난 보편성의 계층을 시사한다:
	\begin{itemize}
		\item \textbf{제1차 (YUCT):} 기본 조정 오차 법칙
		$\varepsilon = \kappa_c \alpha K_{\mathrm{eff}}^{-2/3}$은 가장 깊은 수준에서
		작용하며, 특정 기판에 의존하지 않는다.
		\item \textbf{제2차 (수송):} 2차원 쇼트키 헤테로구조에서 전자 수송에서는,
		운동 지수 $\beta_{\mathrm{IV}}$는 프랙탈 차원 $d_f$로부터
		$\beta_{\mathrm{IV}} = d_f/2$ (가설)로 나타난다.
		\item \textbf{제3차 (성장):} KPZ 부류에 속하는 표면 성장 현상에서는,
		같은 운동 지수가 동적 임계 지수 $z = 3/2$로 나타나며, 성장 계면의 $d_f = 3$의
		성질을 반영한다 (해석).
	\end{itemize}
	따라서 $3/2$ 지수는 재료 특정 우연이 아니라, $d_f = 3$의 조정 과정의
	기하학적 표지이다.
	
	\begin{figure}[htbp]
		\centering
		\begin{tikzpicture}[
			node distance=1.8cm,
			box/.style={rectangle, draw=blue!60, fill=blue!10, rounded corners, 
				text width=3.5cm, align=center, font=\small},
			arrow/.style={->, >=stealth, thick}
			]
			\node[box, fill=red!15, draw=red!50!black] (yuct) 
			{\textbf{제1차 (YUCT)}\\[2pt]
				$\varepsilon = \kappa_c\alpha K_{\mathrm{eff}}^{-2/3}$\\
				(기본 오차 법칙)};
			\node[box, right=of yuct] (transport) 
			{\textbf{제2차 (수송)}\\[2pt]
				$\beta_{\mathrm{IV}} = d_f/2$\\
				(쇼트키 헤테로구조)};
			\node[box, right=of transport] (growth) 
			{\textbf{제3차 (성장)}\\[2pt]
				$z = 3/2$\\
				(KPZ 표면 거칠기)};
			\draw[arrow] (yuct) -- (transport);
			\draw[arrow] (transport) -- (growth);
		\end{tikzpicture}
		\caption{보편성의 계층. YUCT 조정 오차 법칙 (제1차)은 유효 프랙탈 차원 $d_f$와
			전자 수송 (제2차) 및 표면 성장 (제3차)에서 관측되는 운동 지수를 결정한다.
			점선 화살표는 가설을 나타낸다.}
		\label{fig:hierarchy}
	\end{figure}
	
	이 견해는 쇼트키 계면의 프랙탈 차원이 소자 특성에 직접 영향을 미친다는
	증가하는 실험적 증거에 의해 더욱 지지된다~\cite{FractalSchottky}. 이 연구들에서는
	프랙탈 차원이 외부에서 부과된 기하학적 성질 (거칠기)이다. YUCT에서는 이상적인
	측면 접촉에 대한 $d_f = 3$의 차원이 \emph{창발적}이다——그것은 기하학적 거칠기가
	0에 가까워져도 조정 네트워크 자체의 역학으로부터 발생한다. 따라서 외부에서
	부과된 $d_f$는 기존의 내부 조정 네트워크를 변조하며, 조정 기반 소자 설계에
	대한 유망한 경로를 제공한다.
	
	\section{50 자릿수에 걸친 보편성}
	표~\ref{tab:universal-beta}는 다양한 YUCT 부록으로부터, 오차 스케일링 지수가
	$\beta = 2/3$인 것으로 경험적으로 결정되었거나 예측된 시스템들을 수집한다.
	측면 쇼트키 헤테로구조와 KPZ 성장 계면은 추출된 프랙탈 차원 $d_f = 3$에
	기초하여 동일한 지수가 예측되는 시스템으로 추가된다.
	이 표는 관련 에너지 (또는 길이) 스케일에서 50 자릿수 이상을 다룬다.
	
	\begin{table}[H]
		\centering
		\caption{보편적인 YUCT 오차 지수 $\beta = 2/3$ ($0.67$)을 나타내거나 예측되는 시스템.}
		\label{tab:universal-beta}
		\begin{tabular}{@{}llll@{}}
			\toprule
			\textbf{시스템} & \textbf{관측량} & \textbf{상태} & \textbf{참조}\\
			\midrule
			화학 결합 (873)      & $E \propto r^{-0.67}$               & 확인됨  & 부록 C \\
			DNA 복제 오차    & $\varepsilon \propto \Keff^{-0.67}$ & 확인됨  & 부록 E \\
			대사 스케일링 (단순)& $B \propto M^{0.67}$               & 확인됨  & 부록 D, E.7 \\
			GDP 변동 대 태양 활동   & $\sigma_{\mathrm{GDP}} \propto W^{-0.67}$ & 확인됨 & 부록 F \\
			CMB 온도 변동    & $\delta T/T \propto \Keff^{-0.67}$ & 확인됨  & 부록 M \\
			중력파       & $\delta\tau/\tau \propto M^{-0.67}$& 확인됨  & 부록 G, V \\
			\textbf{측면 쇼트키} & $S_I/I^2 \propto T^{1/3}$          & \textbf{예측} & 본 부록 \\
			\textbf{KPZ 계면}\textsuperscript{*} & 폭 변동 $\propto \Keff^{-2/3}$& \textbf{예측} & 본 부록, \cite{KPZ}\\
			\bottomrule
		\end{tabular}
		
		\vspace{0.3cm}
		{\footnotesize \textsuperscript{*}KPZ 성장 계면에 대한 $K_{\mathrm{eff}}$의 정의는 수송의 경우와의 유추로 구축되어야 한다. 예를 들어 상관 부피 내의 상관 자유도의 비율로. 이 예측은 그러한 구축을 가정하며, 테스트되기를 기다린다.}
	\end{table}
	
	\section{결론}
	Ang 등에 의해 보고된 전류-온도 스케일링 지수 $\beta_{\mathrm{IV}} = 3/2$ 및 $1$은
	2차원 헤테로구조에서 열전자 수송을 지배하는 조정 네트워크의 유효 프랙탈 차원의
	운동적 현현이다. YUCT 관계 $\beta_{\mathrm{IV}} = d_f/2$ (가설)와
	$\beta_{\mathrm{err}} = 2/d_f$를 통해, 실험 데이터는 측면 접촉에 대해
	$d_f = 3$을 제공하며, 이는 조정 오차 지수 $\beta_{\mathrm{err}} = 2/3$ ($\approx 0.67$)을
	예측한다. 전류 노이즈 분광법으로 검증 가능한 이 예측은 YUCT 보편 스케일링 법칙의
	적용 범위를 고체 나노전자공학 영역으로 확장한다. 동일한 프랙탈 차원과 지수는
	KPZ 표면 성장에 자연스럽게 나타나며, 비평형 현상에 걸친 깊은 기하학적 통일을
	시사한다. YUCT는 이 통일의 근본적 근거를 제공하지만, 이론과 쇼트키 운동 관계
	및 KPZ 방정식 모두와의 연결은 아직 가설로 남아 있으며, 추가 조사가 정당화된다.
	두 맥락 모두에서 $\Keff$의 운영적 정의는 실험적 테스트를 용이하게 하기 위해
	제공되었다.
	
	\begin{thebibliography}{99}
		\bibitem{Ang2018} Y.~S.~Ang, H.~Y.~Yang, and L.~K.~Ang,
		Phys. Rev. Lett. \textbf{121}, 056802 (2018).
		\bibitem{KPZ} M.~Kardar, G.~Parisi, Y.-C.~Zhang,
		Phys. Rev. Lett. \textbf{56}, 889 (1986).
		\bibitem{FractalSchottky} S.~Zeybek, M.~Biber, A.~Türüt,
		Appl. Surf. Sci. \textbf{253}, 3881 (2007). [쇼트키 다이오드에서 프랙탈 차원 효과에 관한 대표적 연구]
		\bibitem{Yakushev2026} A.~V.~Yakushev, \emph{Yakushev Unified
			Coordination Theory (YUCT)}, Zenodo (2026).
		\bibitem{AppendixL} A.~V.~Yakushev, ``Appendix L: Fractal
		Coordination Error Scaling'', in \emph{YUCT} (2026).
		\bibitem{AppendixM} A.~V.~Yakushev, ``Appendix M: YUCT
		Cosmology'', in \emph{YUCT} (2026).
		\bibitem{AppendixE} A.~V.~Yakushev, ``Appendix E: Coordination
		Genetics'', in \emph{YUCT} (2026).
		\bibitem{AppendixF} A.~V.~Yakushev, ``Appendix F: Application of
		YUCT to Economics'', in \emph{YUCT} (2026).
		\bibitem{AppendixG} A.~V.~Yakushev, ``Appendix G: Resolution of
		the Quantum Gravity Problem'', in \emph{YUCT} (2026).
		\bibitem{AppendixV} A.~V.~Yakushev, ``Appendix V: Time as
		Coordination Entropy'', in \emph{YUCT} (2026).
	\end{thebibliography}
	
\end{document}